\section{The restricted recognizer and integration design}
\label{sec:algorithm}

\subsection{A monotone simplification class}

Fix functions $b(N)$ and $k(N)$, a deterministic RI/RII exhaustion order,
and a deterministic ordering of the birth-front search tree (for example,
iterative depth followed by lexicographic triangle keys).  Let
$\mathsf{BF}_{b,k}$ denote the resulting deterministic preprocessor.  Define
$\mathcal{U}^{\mathsf{BF}}_{b,k}$ to be the set of unknot diagrams on which this
specific process reaches the crossing-free circle:
\begin{enumerate}
\item exhaust RI/RII in the fixed order;
\item if crossings remain, run the complete birth-front search with bounds
      $b(N),k(N)$ and select its first trace in the fixed ordering;
\item apply that trace, return to step 1, and return $\Unknown$ if no trace is
      found.
\end{enumerate}
The policy qualification is necessary.  Two different successful unlock traces
can lead to different reduced diagrams, and the existence of one globally good
sequence does not by itself justify a greedy choice without backtracking.  The
definition permits the diagram size to fall during the process; using the
original $N$ as the upper-bound parameter only weakens the estimate.

\begin{theorem}[Restricted-class unknot recognition]
\label{thm:restricted-recognition}
There is a deterministic, sound algorithm that returns $\Unk$ on every diagram
in $\mathcal{U}^{\mathsf{BF}}_{b,k}$ and otherwise returns either $\Unk$ or $\Unknown$.  Its
running time is
\[
  N^{b(N)+O(1)}(Ck(N))^{k(N)}
\]
and its working memory is polynomial, apart from any subsequent fallback.
\end{theorem}

\begin{proof}
At each stalled state, run the complete birth-front search with the specified
bounds and fixed ordering.  If it succeeds, select the first trace in that
ordering, replay it, expose RI/RII, and exhaust reductions in the fixed order.  Reidemeister moves preserve the represented
knot type.  Therefore reaching the crossing-free circle is a sound unknot
certificate.

Each successful episode is followed by an RI or RII move and hence removes at
least one crossing.  There are at most $N$ episodes.  Multiplying the
single-episode bound from Corollary~\ref{cor:r3-bound} by $N$ is absorbed in the
$N^{O(1)}$ factor.  If a search fails, return $\Unknown$ or invoke an exact
fallback; no negative conclusion is drawn.
\end{proof}

\begin{corollary}
The policy class $\mathcal{U}^{\mathsf{BF}}_{b,k}$ is recognized in quasi-polynomial time when
$b(N),k(N)=O(\log N)$, in $N^{O(\log\log N)}$ time when
$b(N)=O(1)$ and $k(N)=O(\log N)$, and in polynomial time when
$b(N)=O(1)$ and $k(N)=O(\log N/\log\log N)$.
\end{corollary}

\subsection{Certificate size and replay}

A successful preprocessor emits the same kind of trace already used by ProveIt:
for each RIII move, the three stable input crossing indices, followed by RI/RII
records.  Across at most $N$ reduction episodes, a uniform depth bound $k$ gives
$O(Nk)$ RIII records.  Each crossing index uses $O(\log N)$ bits, so the
certificate size is $O(Nk\log N)$ bits.  The existing independent replay routine
checks legality move by move and rebuilds a validated spherical one-component
diagram.

This proof-carrying behavior is more important than the heuristic that found the
trace.  Search policies, budgets, and branch ordering may change without
changing the trusted replay kernel.

\subsection{Opt-in API}

The supplied review patch adds two optional arguments to
\code{simplify}:
\begin{center}
\begin{tabular}{@{}lll@{}}
\toprule
Argument & Suggested values & Meaning \\
\midrule
\code{r3\_search} & \code{"last"}, \code{"clustered"} & legacy or birth-front search \\
\code{r3\_births} & positive integer & maximum number of births \\
\bottomrule
\end{tabular}
\end{center}
The default remains \code{"last"}; therefore applying the patch alone does not
change existing behavior.  In clustered mode, \code{r3\_births=1} gives the
accumulated-support search.  Larger values admit independent birth branches.
The existing \code{r3\_depth}, \code{r3\_budget}, and face-shortcut controls are
retained.

\subsection{Pseudocode}

A simplified version of one depth-limited search is shown below.

\begin{lstlisting}[language=Python,caption={Birth-front RIII search, schematic}]
def search(state, depth_left, births_left, birth_phase,
           active_footprint, last_birth_key, path):
    if budget_exhausted() or depth_left == 0:
        return failure

    if birth_phase and births_left > 0:
        for T in all_legal_triangles(state):
            F = footprint(state, T)
            if F intersects active_footprint:
                continue
            if key(T) <= last_birth_key:
                continue
            apply(T)
            if RI_or_RII_now_legal(F):
                return success
            if search(state, depth_left - 1, births_left - 1,
                      True, active_footprint union F, key(T), path + [T]):
                return success
            undo(T)

    for T in legal_triangles_meeting(state, active_footprint):
        F = footprint(state, T)
        apply(T)
        if RI_or_RII_now_legal(F):
            return success
        if search(state, depth_left - 1, births_left,
                  False, active_footprint union F,
                  last_birth_key, path + [T]):
            return success
        undo(T)

    return failure
\end{lstlisting}

The production helper contains duplicate suppression, inverse-move avoidance,
trial accounting, and the last-move face shortcut.

\subsection{Required production tests}

Before clustered mode can be enabled by default, the following checks are
mandatory.

\begin{enumerate}
\item \textbf{Trace replay.}  Replay every successful trace with the existing
      independent checker.
\item \textbf{Exact restoration.}  Hash the entire $\alpha$ involution before
      and after every failed search and every budget cutoff.
\item \textbf{Invariant preservation.}  Compare Alexander data, writhe, and
      Khovanov ranks on randomized successful RIII traces, as the existing
      synthesis did for the legacy helper.
\item \textbf{Verdict equality.}  Compare the full recognizer with clustered
      mode disabled and enabled on maintained and random corpora.
\item \textbf{Adversarial branching.}  Add fixtures requiring two disjoint
      birth footprints before a connecting RIII move, if a genuine knot-diagram
      example can be found.
\item \textbf{Performance separation.}  Measure successes and failures
      separately.  Deep failed searches should remain budgeted and should run
      only immediately before the expensive exact scan.
\end{enumerate}

\subsection{Why the exact fallback remains essential}

Even a complete $(b,k)$ search is incomplete for unrestricted unknot diagrams.
If it fails, the diagram may still be an unknot that needs larger parameters,
crossing-increasing moves, a different diagram language, or a non-Reidemeister
certificate.  The ProveIt pipeline should therefore keep its exact Khovanov
recognizer and every sound obstruction.  Clustered RIII search is a front-end
optimization with a theorem, not a replacement decision procedure.
